\chapter{ĐÁNH GIÁ VÀ HƯỚNG PHÁT TRIỂN}

\section{Đánh giá Kết quả Đạt được}
Trải qua hai giai đoạn phân tích, thiết kế và hiện thực hóa (Giai đoạn 1 và Giai đoạn 2), nhóm đã hoàn thành toàn diện các mục tiêu đề ra cho bài tập lớn môn Công nghệ Phần mềm:
\begin{itemize}
    \item \textbf{Chuẩn hóa Thiết kế Hệ thống:} Xây dựng hoàn chỉnh bộ tài liệu kỹ thuật từ sơ đồ kiến trúc thành phần, sơ đồ triển khai Docker, sơ đồ lớp Backend, sơ đồ thực thể liên kết (ERD) 12 bảng cho đến chi tiết Use Case, sơ đồ Hoạt động (Activity) và sơ đồ Tuần tự (Sequence) cho toàn bộ 7 phân hệ chức năng.
    \item \textbf{Hiện thực hóa Đầy đủ 7 Phân hệ Nghiệp vụ:} Cài đặt đồng bộ từ giao diện người dùng Next.js 15 đến hệ thống API NestJS 11 và cơ sở dữ liệu PostgreSQL 16. Mọi tương tác giao diện và gọi dịch vụ đều bám sát thiết kế.
    \item \textbf{Làm chủ các Bài toán Kỹ thuật Nâng cao:} Giải quyết trọn vẹn bài toán máy trạng thái đơn hàng (9 chuyển tiếp), bài toán kiểm soát giao dịch thanh toán bất biến (\textit{Idempotency}) và bài toán kiểm soát xung đột tồn kho bằng khóa lạc quan (\textit{Optimistic Locking}).
\end{itemize}

\section{Các Hạn chế Hiện tại}
Bên cạnh những kết quả đã đạt được, hệ thống BrewLite vẫn tồn tại một số điểm giới hạn kỹ thuật cần được ghi nhận trung thực:
\begin{enumerate}
    \item \textbf{Gộp chung mã lỗi xung đột tồn kho và phiên bản:} Tại tầng Backend (\texttt{OrdersService}), cả hai trường hợp hết tồn kho thực tế (\texttt{stock < quantity}) và xung đột phiên bản Khóa lạc quan (\texttt{version} không khớp) hiện đang được xử lý chung thông qua ngoại lệ \texttt{400 BadRequestException}. Việc chưa phân tách riêng mã \texttt{409 Conflict} cho trường hợp xung đột phiên bản khiến phía Client chưa thể kích hoạt cơ chế tự động thử lại (Auto-Retry) ngầm mà người dùng phải thao tác lại thủ công.
    \item \textbf{Phân hệ Loyalty chỉ mới kích hoạt chiều tích điểm (Earn Points):} Phiên bản MVP hiện tại đã hoàn thiện trọn vẹn luồng tự động tích lũy điểm thưởng khi thanh toán thành công (\(1.000 \text{ VNĐ} = 1 \text{ điểm}\)). Tính năng dùng điểm để giảm trừ trực tiếp vào đơn hàng mới (\textit{Redeem Points}) và xử lý hoàn điểm tự động khi hủy đơn được lên kế hoạch bàn giao ở bản phát hành tiếp theo.
    \item \textbf{Giao diện Quầy Pha chế (Barista UI) ngoài phạm vi MVP:} Dù Backend đã hoàn thiện trọn vẹn logic chuyển tiếp trạng thái (\texttt{PREPARING}, \texttt{READY}, \texttt{COMPLETED}), nhưng giao diện web/tablet dành riêng cho nhân viên pha chế tại quầy chưa được hiện thực hóa trong phạm vi bản phát hành MVP này.
    \item \textbf{Tích hợp Cổng thanh toán Thực tế:} Hệ thống hiện đang sử dụng dịch vụ Mock Payment Gateway để giả lập luồng thanh toán. Việc tích hợp các cổng thanh toán thương mại thực tế tại Việt Nam (như VNPay, MoMo, ZaloPay) đòi hỏi quy trình đăng ký doanh nghiệp và thẩm định pháp lý phức tạp nên được định hướng ở giai đoạn sau.
\end{enumerate}

\section{Hướng Phát triển trong Tương lai}
Nhằm đưa BrewLite trở thành một nền tảng vận hành hoàn chỉnh cho các chuỗi cửa hàng cà phê thương mại, các định hướng phát triển tiếp theo bao gồm:
\begin{itemize}
    \item \textbf{Hiện thực hóa tính năng Đổi điểm thưởng (Redeem Loyalty Points):} Cung cấp giao diện cho phép khách hàng chủ động chọn số điểm muốn sử dụng tại bước thanh toán để khấu trừ trực tiếp vào giá trị đơn hàng, tích hợp cơ chế hoàn trả điểm tự động khi đơn hàng bị hủy theo chuẩn cấu trúc cơ sở dữ liệu đã thiết kế sẵn.
    \item \textbf{Xây dựng Màn hình Pha chế (Barista POS Display):} Phát triển giao diện chuyên dụng cho quầy bar, cập nhật danh sách đơn hàng theo thời gian thực sử dụng công nghệ WebSockets hoặc Server-Sent Events (SSE).
    \item \textbf{Tích hợp Cổng thanh toán VNPay / MoMo:} Kết nối trực tiếp qua Webhook và chuẩn thanh toán mã VietQR, hỗ trợ hoàn tiền tự động khi đơn hàng bị hủy sau khi đã thanh toán.
    \item \textbf{Phân tách mã lỗi 409 Conflict và Cơ chế Client-side Auto-Retry:} Nâng cấp \texttt{OrdersService} để phân tách rõ ràng giữa mã lỗi \texttt{400 Bad Request} (khi hết hàng thực tế) và mã lỗi \texttt{409 Conflict} (khi tồn kho còn nhưng xảy ra xung đột version giữa các request đồng thời), kết hợp xây dựng lớp client tự động thử lại (Retry với Exponential Backoff) giúp nâng cao trải nghiệm người dùng trong các đợt cao điểm.
    \item \textbf{Phân hệ Báo cáo và Quản trị Doanh thu (Admin Analytics):} Bổ sung biểu đồ trực quan hóa dữ liệu bán hàng, dự báo nhu cầu nguyên vật liệu và quản lý các chiến dịch khuyến mãi theo mùa.
\end{itemize}
